\scanpage{188}
\mspage{393}

2) \uline{Determination by} $\alpha_0,\beta_0,f$.

There is a one-to-one correspondence between the pairs $(g,h)$ satisfying (25), and the pairs
\[
\tag{31}
(\alpha_0,\beta_0)\qquad
\bigl(\alpha_0\in\Pi_{0,3}^{\wedge\sim}=\Pi_{0,3}^{\wedge}\cdot\{1,\tau_\infty\},\quad
\beta_0\in\Pi_{0,3}^{\wedge}\bigr)
\]
\deleted{\unsure{satisf}} via the relations
\[
\tag{32}
\left\{\begin{aligned}
g&=\alpha_0\rho(\alpha_0)^{-1}=[\alpha_0,\rho]\\
h&=\beta_0\sigma_\infty(\beta_0)^{-1}=[\beta_0,\sigma_\infty]
\end{aligned}\right.
\]
Thus, the $u\in\unsure{\mathcal{M}_{0,3}^{!\sim}|\mathcal{N}_{0,3}^{!\sim}}$ correspond to the systems
\[
\tag{33}
(\alpha_0,\beta_0,f)\qquad
\alpha_0\in\Pi_{0,3}^{\sim\wedge},\quad\beta_0,f\in\Pi_{0,3}^{\wedge},\quad
f\ \text{\unsure{determined} mod }f\mapsto\ell_0^n f\ (n\in\widehat{\mathbb Z})
\]
satisfying the condition \inserted{on loops} (29), considered as a condition in $\alpha_0,\beta_0,f$, \deleted{\unsure{by}} substituting in it \inserted{for} $g,h$ the values given by (32). The relation is written
\[
\sigma_\infty(\beta_0)\beta_0^{-1}\alpha_0\rho(\alpha_0)^{-1}
=\sigma_\infty(f^{-1})\ell_1^\nu\rho(f)
\]
or again
\[
\sigma_\infty(f\beta_0)
\underbrace{\beta_0^{-1}\alpha_0}_{(f\beta_0)^{-1}(f\alpha_0)}
\rho(\inserted{f}\alpha_0\deleted{\unsure{f}})^{-1}
=\ell_1^\nu.
\]
or again
\[
\tag{34}
\boxed{(f\alpha_0)\rho(f\alpha_0)^{-1}
=\bigl((f\beta_0)\sigma_\infty(f\beta_0)^{-1}\bigr)\ell_1^\nu}
\]
(loop relation in $\alpha_0,\beta_0,f$). As with (29), $\nu\in\widehat{\mathbb Z}$ is uniquely determined in terms \unsure{of|of the} $\alpha_0,\beta_0$, by the possibility of finding $f$ satisfying relation (34).

The advantage of introducing $\alpha_0,\beta_0$ instead

\scanpage{189}

of $g,h$ is to make unnecessary (since \unsure{automatic}) the relations (25), -- only relation (34) remains to be checked. The relations (23) describing $u$ become here:
\[
\tag{35}
u(\sigma_\infty)=\operatorname{int}(\beta_0)(\sigma_\infty),\qquad
u(\rho)=\operatorname{int}(\alpha_0)(\rho)
\]

3) \uline{Determination by} $(\alpha,\beta,f)$.

The particular form of (34), where the triple $(\alpha_0\ \beta_0\ f)$ enters only through the corresponding pair $(f\alpha_0,f\beta_0)$, leads to setting
\[
\tag{35}
\alpha=f\alpha_0\deleted{\unsure{=}},\qquad\beta=f\beta_0,
\]
we find that the $u\in\mathcal{M}_{0,3}^{!\sim}$ are determined by the triples
\[
\tag{36}
(\alpha,\beta,f)\qquad
\begin{gathered}
\alpha\in\Pi_{0,3}^{\sim\wedge},\quad\beta,f\in\Pi_{0,3}^{\wedge},\quad
\alpha,\beta,f\ \text{modulo}\\
\alpha\mapsto\ell_0^n\alpha,\quad\beta\mapsto\ell_0^n\beta,\quad
f\mapsto\ell_0^n f\ (n\in\widehat{\mathbb Z})\\
\text{(the \uline{same} }n\text{ for }\alpha,\beta,f\text{)}
\end{gathered}
\]
\margin{A disadvantage is the indeterminacy of the pair $(\alpha,\beta)$, under $(\alpha,\beta)\longrightarrow(\ell_0^n\alpha,\ell_0^n\beta)$ $(n\in\widehat{\mathbb Z})$, but this is a minor drawback, which \unsure{will be removed} below.}
satisfying the condition
\[
\tag{37}
\boxed{\bigl(\alpha\rho(\alpha)^{-1}\bigr)
=\bigl(\beta\sigma_\infty(\beta)^{-1}\bigr)\cdot\ell_1^\nu}
\]
(\uline{loop relation in} $\alpha,\beta$). The advantage of this relation is that $f$ no longer appears in it.

We recover $\alpha_0,\beta_0,g,h$ in terms of $\alpha,\beta,f$ by the formulas
\[
\tag{38}
\alpha_0=f^{-1}\alpha,\qquad\beta_0=f^{-1}\beta
\]
\[
\tag{39}
g=\alpha_0\rho(\alpha_0)^{-1}
=f^{-1}\bigl(\underbrace{\alpha\rho(\alpha)^{-1}}_{\deleted{\illegible{1}}}\bigr)\rho(f)
=f^{-1}\xi\rho(f)
\]
\[
\tag{40}
h=\beta_0\sigma_\infty(\beta_0)^{-1}
=f^{-1}\bigl(\beta\sigma_\infty(\beta)^{-1}\bigr)
=f^{-1}\eta\sigma_\infty(f),
\]
where we have set
\[
\tag{41}
\xi=\alpha\rho(\alpha)^{-1},\qquad
\eta=\beta\sigma_\infty(\beta)^{-1}
\qquad(\xi,\eta\in\Pi_{0,3}^{\wedge})
\]

\scanpage{190}
\mspage{394}

The formulas (23) -- (35) determining $u$ in terms of the \unsure{computational} data, become here
\[
\tag{42}
u(\sigma_\infty)=\operatorname{int}(f^{-1}\beta)(\sigma_\infty)
\qquad u(\rho)=\operatorname{int}(f^{-1}\alpha)(\rho)
\]
\inserted{i.e.
\[
\tag{42bis}
\text{\deleted{\unsure{i.e.}}}\ u=\operatorname{int}(f^{-1})u_{\alpha,\beta},
\]
where $u_{\alpha,\beta}$ depends only on $\alpha,\beta$:
\[
\tag{42ter}
\boxed{u_{\alpha,\beta}(\sigma_\infty)=\operatorname{int}(\beta)(\sigma_\infty)=\unsure{\eta\sigma_\infty}}
\quad
\boxed{u_{\alpha,\beta}(\rho)=\operatorname{int}(\alpha)(\rho)=\xi\rho}
\]}
Formula (37) \unsure{can be interpreted as a} relation in $\xi,\eta$
\[
\tag{43}
\xi=\eta\ell_1^\nu.
\]

4) \uline{Determination by} $(\xi,\eta,f)$.

The datum of the pair $(\alpha,\beta)\in\Pi_{0,3}^{\overset{\inserted{\wedge}}{\sim}\wedge}\times\Pi_{0,3}^{\wedge}$, is equivalent to that of the pair $(\xi,\eta)$ deduced from it by (41), \unsure{subject} to the \deleted{formulas}\inserted{equations}
\[
\tag{44}
\boxed{\left\{\begin{aligned}
\xi\rho(\xi)\rho^2(\xi)&=1\\
\eta\sigma_\infty(\eta)&=1
\end{aligned}\right.}
\]
Thus, the datum of $u\in\mathcal{M}_{0,3}^{!\sim}$ is equivalent to that of a triple
\[
\tag{45}
(\xi,\eta,f)\qquad
\begin{gathered}
\xi,\eta,f\in\Pi_{0,3}^{\wedge}\quad\text{modulo}\\
\xi\mapsto\ell_0^n\xi\unsure{\ell_1^n|\ell_1^{-n}},\quad
\eta\mapsto\ell_0^n\eta\unsure{\ell_1^n|\ell_1^{-n}},\quad
f\mapsto\ell_0^n f
\end{gathered}
\]
satisfying the relations (44), and the relation
\[
\tag*{(46)=(43)}
\boxed{\xi=\eta\ell_1^\nu},
\]
which plays the role of ``\uline{loop relation in} $\xi,\eta$''. The advantage of this presentation is

\scanpage{191}

the maximal simplicity of the equations (44),(46) -- and the fact that $f$ does not enter into them: \deleted{\unsure{these are}} the first two equations are equations in $\xi$ alone and in $\eta$ alone, and the third relates $\xi$ and $\eta$ in a particularly simple way -- in fact if $\mu=1$, i.e. $\nu=0$, it is the relation $\xi=\eta$ !

The drawback is that there are three equations instead of one, and that we have no explicit expression for $u(\sigma_\infty)$ and $u(\rho)$ in terms of $\xi,\eta,f$: we must first, for this, solve the equations
\margin{but yes, we do have a nice explicit \unsure{expression} (42ter):
\[
u(\sigma_\infty)=\operatorname{int}(f^{-1})(\eta\sigma_\infty)
\]
\[
u(\rho)=\operatorname{int}(f^{-1})(\xi\rho)
\]
!!}
\[
\tag{47}
\left\{\begin{aligned}
\alpha\rho(\alpha)^{-1}&=\xi\qquad(\alpha\in\Pi_{0,3}^{\sim\wedge})\\
\beta\sigma_\infty(\beta)^{-1}&=\eta\qquad(\beta\in\Pi_{0,3}^{\wedge})
\end{aligned}\right.
\]
in $\alpha,\beta$, to be able then to apply (42).

Another \inserted{[!]} drawback, of importance, is that when we work with the images of the elements under consideration in $\mathrm{Gl}(2,\widehat{\mathbb Z})^{\unsure{\pm\prime|\pm1}}$, the knowledge in this group of $\xi,\eta$ \inserted{(even in $\mathrm{Sl}(2,\widehat{\mathbb Z})$ itself)} is \deleted{\unsure{too}} a rather weak datum, which in no way allows us to recover $\alpha,\beta$ as elements of $\mathrm{Gl}(2,\widehat{\mathbb Z})^{\unsure{\pm\prime|\pm1}}$, cf. below.

\scanpage{192}
\mspage{395}

5) \uline{Composition \unsure{with the|with} automorphisms} \uline{that are inner}. \uline{Cocycle relations} \uline{for a composite} $u''=u'u$.

Let $\gamma\in\Pi_{0,3}^{\wedge}$, and let us consider the composite
\[
\tag{48}
u'=\operatorname{int}(\gamma)\circ u
\]
it corresponds to invariants
\[
f',g',\unsure{h'|\beta'},\alpha_0',\beta_0',\alpha',\beta',\xi',\eta'
\]
given by
\[
\tag{49}
\left\{\begin{aligned}
f'&=f\gamma^{-1},\quad\deleted{\unsure{g'=\gamma g},\ h'=\gamma h}\\
g'&=\gamma g\rho(\gamma)^{-1},\quad h'=\gamma h\sigma_\infty(\gamma)^{-1}\\
\alpha_0'&=\gamma\alpha_0,\quad\beta_0'=\gamma\beta_0\\
\alpha'&=\alpha,\quad\beta'=\beta\\
\xi'&=\xi,\quad\eta'=\eta
\end{aligned}\right.
\]
Thus $\alpha$ and $\beta$ do not change when we replace $\unsure{u|f}$ by $\operatorname{int}(\gamma)u$, they depend only on the outer loop automorphism defined by $u$, i.e. on the image of $u$ in $\mathcal{N}_{0,3}^{!\sim}$. This is a substantial advantage of the use of $(\alpha,\beta,f)$ or of $(\xi,\eta,f)$, in preference to $(\alpha_0,\beta_0,f)$ or of $(g,h,f)$.

Now let
\[
\tag{50}
u,u'\in\mathcal{M}_{0,3}^{!\sim},\quad\text{let }u''=u'u
\]
let
\[
\tag{51}
\left\{\begin{aligned}
&f,g,h,\alpha_0,\beta_0,\alpha,\beta,\xi,\eta\\
&f',g',h',\alpha_0',\beta_0',\alpha',\beta',\xi',\eta'\qquad\unsure{\text{and}}\\
&f'',g'',h'',\alpha_0'',\beta_0'',\alpha'',\beta'',\xi'',\eta''
\end{aligned}\right.
\]

\scanpage{193}

be the invariants relative to $u,u',u''$ respectively, -- we \unsure{want} to express (as far as possible) those of $u''$ in terms of those of $u,u'$, and we find
\[
\begin{aligned}
&\left\{\begin{array}{lll}
g''=u'(g)g' & \alpha_0''=u'(\alpha_0)\alpha_0'\, {*} & \alpha''=\unsure{(f'u')}(\alpha).\alpha'\, {*}\\
h''=u'(h)h' & \beta_0''=u'(\beta_0)\beta_0' & \beta''=(f'u')(\beta).\beta'
\end{array}\right.\\
&\text{but } f''=f'u'(f)\deleted{\unsure{\xi}}.
\end{aligned}
\tag{52}
\]
\margin{${*}$ $u'(\alpha_0)$ and $(f'u')(\alpha)$ make sense a priori only if $\alpha_0,\alpha\in\widehat{\Pi}_{0,3}$ i.e. $\mu\equiv1\ (3)$, otherwise (for example) we can still define $u'(\alpha_0)=u'\alpha_0u'^{-1}$, but it will no longer be \deleted{\unsure{in}} \illegible{2} in $\widehat{\Pi}_{0,3}$, and the formula $\alpha_0''=u'(\alpha_0)\alpha_0'$ will make sense, but will be \uline{false} !}

We thus see that if we consider \unsure{(for} $u$ varying in $\mathcal M_{0,3}^{!\sim})$ $f,g,h,\alpha_0,\beta_0,\alpha,\beta$ as functions
\[
\mathcal M_{0,3}^{!\sim}\longrightarrow\widehat{\Pi}_{0,3},\ \text{or }\widetilde{\Pi}_{0,3}^{\wedge}
\]
then $f,g^{-1},h^{-1},\alpha_0^{-1},\beta_0^{-1}$ satisfy the relation for \uline{cocycles}, and $\alpha^{-1},\beta^{-1}$ do too when we restrict ourselves to the\deleted{\illegible{1}} subgroup for which $f=1$, i.e. the automorphisms $u\in\mathcal M_{0,3}^{!\sim}$ which \uline{normalize} \deleted{\unsure{the}} $L_0=\{\ell_0^n\mid n\in\widehat{\mathbb Z}\}$. To tell the truth, this statement \inserted{\deleted{\illegible{1}} for $f$, $\alpha$, $\beta$} is a little lame, because of the indeterminacies for $f$, entailing a corresponding indeterminacy for $\alpha,\beta$, $f\mapsto\ell_0^n f$, $\alpha\mapsto\ell_0^n\alpha$, $\beta\mapsto\ell_0^n\beta$. To construct canonical cocycles corresponding to $f,\alpha^{-1},\beta^{-1}$, it woul\inserted{d} be necessary to introduce a group extension of $\mathcal M_{0,3}^{!\sim}$ by $\widehat{\mathbb Z}$, on the model of the special Teichmüller groups, on which $f,\alpha,\beta$ would be defined without ambiguity. There is no need here to go into detail\ldots